\ncdlesson{1}{Undskyld, er det plads 14?}{请问，这是 14 号座吗？}
{DU3 Modul 1（A1 起步）｜对标 NCE 1 Lesson 1--4}
{主格人称代词 · 是非问句倒装（动词 + 主语）· \textit{Ja / Nej / Jo} 应答法则 · 简短回答（\textit{Det er det}）}
{第 1 册旨在筑牢发音和语序的“肌肉记忆”。本课不要死背单词表，请重点对着\textbf{语音与音标标注}大声朗读课文 5 遍，直到把 \textit{er det}、\textit{undskyld}、\textit{sidder} 的嘴型练顺，然后一口气做完第五部分的句型改写。}

% ---------------------------------------------------------------------
\begin{foerduleaeser}
在阅读下面的对话前，先记住这两个问题，带着问题读课文：
\begin{enumerate}
  \item \textit{Sidder manden på den rigtige plads i toget?}（那位男士坐对列车座位了吗？）
  \item \textit{Hvor kommer Emma fra, og hvor skal hun stå af?}（艾玛来自哪里，她要在哪一站下车？）
\end{enumerate}
\end{foerduleaeser}

% ---------------------------------------------------------------------
\ncdsubsec{二、Dialog \& Udtale（课文对话、发音标注与中文译文）}

\begin{tipbox}{场景说明}
从哥本哈根开往奥胡斯的 DSB 城际列车上，Emma 拿着车票核对座位号，发现座位上已经坐着一位丹麦男生 Jonas。
\end{tipbox}

\dialogturn{1}{Emma}{Undskyld, er det plads 14?}
{\textit{undskyld} \ipa{ˈɔnsgylˀ}（末尾 \textbf{d 完全不发音}，u 读接近 o，末尾带轻微喉塞音 \ipa{ˀ}）；\textit{er det} 连读弱化为 \ipa{æɐ̯ dɛ}；\textit{plads} \ipa{ˈplas}（\textbf{d 不发音}）。}
{打扰一下，这是 14 号座吗？}

\dialogturn{2}{Jonas}{Ja, det er det. Har du også billet til nummer 14?}
{\textit{det er det} \ipa{de ˈæɐ̯ de}（\textit{det} 末尾的 \textbf{t 永远不发音}！重音落在第一个 \textit{er} 上）；\textit{også} \ipa{ˈʌsə}（\textbf{g 不发音}）；\textit{billet} \ipa{biˈlɛd}（法语借词，重音在第二音节）。}
{是的，是 14 号。你也有 14 号的票吗？}

\dialogturn{3}{Emma}{Ja, jeg tror, du sidder på min plads. Se: tog 23, vogn 7, plads 14.}
{\textit{jeg} \ipa{ˈjɑɪ̯}；\textit{sidder} \ipa{ˈseðɐ}（\textit{dd} 读作\textbf{软 d \ipa{ð}}——舌尖轻抵下齿背，气流从舌两侧轻擦出来，不要读成英文的 d；词尾 \textit{-er} 弱读成 \ipa{ɐ}）；\textit{tog} \ipa{ˈtɔʊ̯ˀ}；\textit{vogn} \ipa{ˈvʌʊ̯ˀn}。}
{是的，我想你坐了我的座位。你看：23 次列车，7 号车厢，14 号座。}

\dialogturn{4}{Jonas}{Åh, undskyld! Er det ikke vogn 6?}
{\textit{ikke} \ipa{ˈeɡə}，但在日常语流里常被吞成短促的 \ipa{eg}；\textit{er det ikke} 连读听起来像 \ipa{æɐ̯ dɛ eg}。}
{哦，真抱歉！难道这不是 6 号车厢吗？}

\dialogturn{5}{Emma}{Nej, det er det ikke. Det er vogn 7. Det er ikke den rigtige vogn for dig.}
{\textit{rigtige} \ipa{ˈʁɛgdiə}（小舌音 r 起头，中间 g 弱化，词尾 -e 轻读 \ipa{ə}）；\textit{dig} \ipa{ˈdɑɪ̯}（词尾 g 滑成 i 音）。}
{不，不是的。这是 7 号车厢。这不是你要坐的车厢。}

\dialogturn{6}{Jonas}{Hvor dumt! Jeg sidder i den forkerte vogn. Er pladsen ved siden af ikke ledig? Ellers går jeg ind i vogn 6.}
{\textit{hvor} \ipa{ˈvɒːˀ}（\textbf{\textit{hv-} 开头的词，字母 h 永远不发音}！）；\textit{dumt} \ipa{ˈdɔmˀd}；\textit{forkerte} \ipa{fɒˈkʰeɐ̯ˀdə}；\textit{ved siden af} \ipa{ve ˈsiːðən ˈæːˀ}（\textit{af} 在此 f 脱落）；\textit{ledig} \ipa{ˈleːði}（\textbf{词尾 -ig 中的 g 不发音}，d 读软 d）。}
{真糊涂！我坐错车厢了。旁边这个座位难道不是空的吗？如果不是的话我就去 6 号车厢。}

\dialogturn{7}{Emma}{Jo, det er den! Det gør ikke noget. Jeg hedder Emma.}
{\textit{Det gør ikke noget} \ipa{de ˈgɶɐ̯ egə ˈnɔːð}（整句是丹麦人嘴边的高频固定句，\textit{noget} 的 t 不发音、g 极弱，听起来接近 \ipa{nɔːð}）；\textit{hedder} \ipa{ˈheðɐ}（软 d）。}
{是空的呀！（这）没关系。我叫艾玛。}

\dialogturn{8}{Jonas}{Tak, hvor er du sød! Jeg hedder Jonas. Kommer du fra København?}
{\textit{sød} \ipa{ˈsøːðˀ}（长圆唇音 \ipa{ø} + 喉塞音 + 软 d）；\textit{København} \ipa{kʰøbənˈhɑʊ̯ˀn}（重音在最后一个音节 \textit{-havn}）。}
{谢谢，你人真好！我叫约纳斯。你是哥本哈根人吗？}

\dialogturn{9}{Emma}{Nej, jeg kommer fra Kina, men jeg bor i København nu. Jeg skal stå af på Aarhus Station.}
{\textit{skal} \ipa{ˈsgal}（句中常弱读为 \ipa{sga}，\textbf{l 经常被吞掉}！）；\textit{stå af} \ipa{sdɔ ˈæːˀ}（注意 \textit{af} 读 \ipa{æːˀ}，不要读出 f）；\textit{station} \ipa{sdaˈɕoːˀn}（\textit{ti} 读接近汉语“西”的 \ipa{ɕ}）。}
{不是，我来自中国，但我现在住在哥本哈根。我要在奥胡斯站下车。}

\dialogturn{10}{Jonas}{Hvor hyggeligt! Jeg skal også stå af i Aarhus. God tur!}
{\textit{hyggeligt} \ipa{ˈhyɡəli}（词尾 \textit{-igt} 中的 \textbf{g 和 t 都不发音}！只读 \ipa{i}）；\textit{God tur!} \ipa{ˈɡoːˀ ˈtʰuɐ̯ˀ}（\textit{god} 的 d 极轻或近乎弱化）。}
{真巧/真棒！我也要在奥胡斯下车。祝你旅途愉快！}

% ---------------------------------------------------------------------
\ncdsubsec{三、Gloser（核心生词表）}

\begin{tipbox}{记词铁律}
丹麦语名词没有性数规律可猜，动词时态有多种变化。从第 1 课起，请连同括号内的屈折形式一起记忆。
\end{tipbox}

\begin{glosertable}
名词 & en plads & \textit{(pladsen, pladser, pladserne)} & 座位；广场；空间 \\
名词 & et tog & \textit{(toget, tog, togene)} & 火车（单复数同形） \\
名词 & en vogn & \textit{(vognen, vogne, vognene)} & 车厢；推车 \\
名词 & en billet & \textit{(billetten, billetter, billetterne)} & 车票；门票 \\
名词 & et nummer & \textit{(nummeret, numre, numrene)} & 号码；数字 \\
名词 & en station & \textit{(stationen, stationer, stationerne)} & 车站 \\
\midrule
动词 & at være & \textit{(er, var, har været)} & 是；在 \\
动词 & at sidde & \textit{(sidder, sad, har siddet)} & 坐着（静态位置） \\
动词 & at stå & \textit{(står, stod, har stået)} & 站着；\textit{stå af} 下车 \\
动词 & at hedde & \textit{(hedder, hed, har heddet)} & 叫……名字 \\
动词 & at komme & \textit{(kommer, kom, er kommet)} & 来；来自（完成时用 \textit{er}） \\
动词 & at tro & \textit{(tror, troede, har troet)} & 认为；相信 \\
\midrule
形容词 & rigtig & \textit{(rigtigt, rigtige)} & 正确的；真正的 \\
形容词 & ledig & \textit{(ledigt, ledige)} & 空闲的；空着的 \\
形容词 & forkert & \textit{(forkert, forkerte)} & 错误的（中性不加 \textit{-t}） \\
形容词 & sød & \textit{(sødt, søde)} & 友好的；可爱的；甜的 \\
\midrule
副词 & ellers & --- & 否则；不然的话；另外 \\
副/叹词 & også / jo & --- & 也 /（反驳否定问句的）是呀 \\
\end{glosertable}

% ---------------------------------------------------------------------
\ncdsubsec{四、Grammatik \& Faste vendinger（核心语法与地道表达精讲）}

\begin{grammatikbox}{语法点 1：主格人称代词与系动词 \textit{er}}
丹麦语动词\textbf{没有人称和数的变化}！不管是“我、你、他”还是“我们、他们”，现在时永远是同一个词形：
\begin{itemize}
  \item 单数：\da{jeg er}（我是）/ \da{du er}（你是）/ \da{han er}（他是）/ \da{hun er}（她是）/ \da{det er}（它是）/ \da{den er}（共性“它”是）
  \item 复数与敬称：\da{vi er}（我们是）/ \da{I er}（你们是，注意 \textbf{I 永远大写}）/ \da{de er}（他们是）/ \da{De er}（您是，敬称）
\end{itemize}
\end{grammatikbox}

\begin{grammatikbox}{语法点 2：是非问句倒装（Ja/nej-spørgsmål）}
英语需要借助助动词 \textit{Do/Does/Is}，而丹麦语\textbf{所有动词直接提前到句首}，形成“\textbf{动词 + 主语}”倒装：

\vspace{0.3em}
\renewcommand{\arraystretch}{1.25}
\begin{tabularx}{\linewidth}{@{} X X @{}}
\toprule
\textbf{陈述句（主语 + 动词 + 其他）} & \textbf{是非问句（\textcolor{NCDRed}{动词 + 主语} + 其他）} \\
\midrule
\textit{Det \textbf{er} plads 14.}（这是 14 号座。） & \textit{\dabold{Er} \da{det} plads 14?}（这是 14 号座吗？） \\
\textit{Du \textbf{kommer} fra Kina.}（你来自中国。） & \textit{\dabold{Kommer} \da{du} fra Kina?}（你来自中国吗？） \\
\textit{Du \textbf{hedder} Jonas.}（你叫约纳斯。） & \textit{\dabold{Hedder} \da{du} Jonas?}（你叫约纳斯吗？） \\
\textit{Du \textbf{skal} til Aarhus.}（你要去奥胡斯。） & \textit{\dabold{Skal} \da{du} til Aarhus?}（你要去奥胡斯吗？） \\
\bottomrule
\end{tabularx}
\vspace{0.3em}

\textbf{注意否定疑问句}：把 \textit{ikke} 放在\textbf{“动词 + 主语”之后}，如：\textit{Er det \dabold{ikke} vogn 6?} / \textit{Kommer du \dabold{ikke} fra København?}
\end{grammatikbox}

\begin{grammatikbox}{语法点 3：\textit{Ja}、\textit{Nej} 与 \textit{Jo} 的黄金法则（考场高频失分点！）}
丹麦语中的“是 / 不是”不是看“客观事实是真是假”，而是看\textbf{问句里有没有否定词（如 \textit{ikke}、\textit{aldrig}、\textit{ingen}）}：

\vspace{0.3em}
\renewcommand{\arraystretch}{1.25}
\begin{tabularx}{\linewidth}{@{} >{\raggedright\arraybackslash}p{4.6cm} X X @{}}
\toprule
\textbf{问句类型} & \textbf{你的真实意思为“肯定（是）”} & \textbf{你的真实意思为“否定（不是）”} \\
\midrule
\textbf{普通肯定问句}（不含 \textit{ikke}）\newline 例：\textit{Kommer du fra Danmark?} & \dabold{Ja}, det gør jeg.\newline （是的，我来自丹麦。） & \dabold{Nej}, det gør jeg ikke.\newline （不是，我不来自丹麦。） \\
\textbf{否定问句}（含 \textit{ikke}）\newline 例：\textit{Er det \textbf{ikke} vogn 6?} & \dabold{Jo}, det er det!\newline （不对，\textbf{是} 6 号车厢！）$\leftarrow$ \textbf{反驳对方} & \dabold{Nej}, det er det ikke.\newline （对，\textbf{确实不是} 6 号车厢。）$\leftarrow$ \textbf{赞同对方} \\
\bottomrule
\end{tabularx}
\vspace{0.3em}

\textbf{一句话口诀}：只要对方在问句里带了 \textit{ikke}，而你想说“\textbf{不对，其实是这样！}”，就\textbf{必须用 \dabold{Jo}}，绝对不能用 \textit{Ja}！课文中第 4--5 句是赞同对方（确实不是 6 号车厢 $\rightarrow$ \textit{Nej}）；第 6--7 句是反驳对方“难道不是空的吗？”（其实是空的！$\rightarrow$ \textit{Jo}）。
\end{grammatikbox}

\begin{grammatikbox}{语法点 4：地道简短回答（Korte svar）与高频固定搭配}
丹麦人在回答是非问句时，很少只干巴巴说一个 \textit{Ja} 或 \textit{Nej}，通常会带上一个简短回指结构：
\begin{itemize}
  \item \textbf{问句动词是 \textit{er / kan / skal / vil / har}}：直接重复该动词 + \textit{det}（或对应代词）。\\
    例：\textit{Er det plads 14?} $\rightarrow$ \da{Ja, det er det.} / \da{Nej, det er det ikke.}\\
    例：\textit{Er pladsen ledig?} $\rightarrow$ \da{Ja, det er den.}（因为 \textit{en plads} 是 \textit{en} 词，可用 \textit{den} 回指）\\
    例：\textit{Skal du til Aarhus?} $\rightarrow$ \da{Ja, det skal jeg.}
  \item \textbf{问句动词是普通实义动词}（如 \textit{kommer, bor, sidder, hedder}）：用万能代动词 \dabold{gør} 替代！\\
    例：\textit{Kommer du fra Kina?} $\rightarrow$ \da{Ja, det gør jeg.} / \da{Nej, det gør jeg ikke.}
\end{itemize}

\textbf{高频地道搭配（Faste vendinger）：}
\begin{itemize}
  \item \da{Det gør ikke noget.} —— 没关系，不要紧。（回应 \textit{Undskyld} 的最常用语）
  \item \da{stå på / stå af (toget / bussen)} —— 上车 / 下车。（乘公交火车上下车固定搭配）
  \item \da{Hvor er du sød! / Hvor dumt!} —— 你真好！/ 真糊涂！（感叹句用 \textit{Hvor + 形容词} 开头）
\end{itemize}
\end{grammatikbox}

% ---------------------------------------------------------------------
\ncdsubsec{五、Mønsterøvelser（新概念式句型操练）}

\textbf{练习 A：将下列陈述句改写为 (1) 普通是非问句，(2) 带 \textit{ikke} 的否定疑问句}
\begin{quote}
\small 例：\textit{Det er vogn 7.} $\rightarrow$ (1) \textit{Er det vogn 7?} \quad (2) \textit{Er det ikke vogn 7?}
\end{quote}
\begin{enumerate}
  \item \textit{Pladsen er ledig.}
  \item \textit{Du bor i København.}
  \item \textit{Jonas sidder i den forkerte vogn.}
  \item \textit{Toget stopper på Odense Station.}
  \item \textit{Du skal stå af i Aarhus.}
\end{enumerate}

\textbf{练习 B：用 \textit{Ja}、\textit{Nej} 或 \textit{Jo} 搭配简短回答完成对话}
\begin{enumerate}
  \item A: \textit{Er det din billet?}（是的，是我的票。）$\rightarrow$ B: \underline{\hspace{1.2cm}}, \underline{\hspace{3.2cm}}.
  \item A: \textit{Bor du ikke i Danmark?}（不对，我确实住在丹麦！）$\rightarrow$ B: \underline{\hspace{1.2cm}}, \underline{\hspace{3.2cm}}.
  \item A: \textit{Kommer Emma fra Danmark?}（不，她不来自丹麦。）$\rightarrow$ B: \underline{\hspace{1.2cm}}, \underline{\hspace{3.2cm}}.
  \item A: \textit{Er det ikke den rigtige station?}（是的，确实不是那一站。）$\rightarrow$ B: \underline{\hspace{1.2cm}}, \underline{\hspace{3.2cm}}.
  \item A: \textit{Skal I stå af her?}（是的，我们要在这里下车。）$\rightarrow$ B: \underline{\hspace{1.2cm}}, \underline{\hspace{3.2cm}}.
\end{enumerate}

\textbf{练习 C：中译丹实战（注意语序与名词冠词）}
\begin{enumerate}
  \item 打扰一下，这个座位是空的吗？
  \item 我想你坐了错误的车厢。（提示：用 \textit{Jeg tror, du...}）
  \item 没关系，你人真好！
  \item 难道你不是在奥胡斯下车吗？（用 \textit{ikke} 提问）
  \item 是的，我要在那里下车。（反驳上面的否定问句，用 \textit{Jo}！）
\end{enumerate}

% ---------------------------------------------------------------------
\ncdsubsec{六、Resumé \& Mundtlig træning（口语复述与参考答案）}

\textbf{口语复述任务（对标 Modul 1--2 口试表达）}：请遮住课文，先用丹麦语口头回答 \textbf{Før du læser} 的两个问题，然后尝试用 \textbf{4--5 句话} 以第三人称复述今天的故事。

\begin{answerbox}
\textbf{练习 A 答案}
\begin{enumerate}
  \item (1) \textit{Er pladsen ledig?} \quad (2) \textit{Er pladsen ikke ledig?}
  \item (1) \textit{Bor du i København?} \quad (2) \textit{Bor du ikke i København?}
  \item (1) \textit{Sidder Jonas i den forkerte vogn?} \quad (2) \textit{Sidder Jonas ikke i den forkerte vogn?}
  \item (1) \textit{Stopper toget på Odense Station?} \quad (2) \textit{Stopper toget ikke på Odense Station?}
  \item (1) \textit{Skal du stå af i Aarhus?} \quad (2) \textit{Skal du ikke stå af i Aarhus?}
\end{enumerate}

\textbf{练习 B 答案}
\begin{enumerate}
  \item \dabold{Ja, det er det.}（肯定问句 $\rightarrow$ 肯定回答用 \textit{Ja}）
  \item \dabold{Jo, det gør jeg!}（否定问句 $\rightarrow$ 反驳否定用 \textit{Jo}；\textit{bor} 是实义动词，简答用 \textit{gør}）
  \item \dabold{Nej, det gør hun ikke.}（肯定问句 $\rightarrow$ 否定回答用 \textit{Nej}；\textit{kommer} 用 \textit{gør} 替代）
  \item \dabold{Nej, det er det ikke.}（否定问句 $\rightarrow$ 赞同对方的否定“确实不是”，用 \textit{Nej}！）
  \item \dabold{Ja, det skal vi.}（\textit{I} 提问“你们”，回答用 \textit{vi}“我们”，重复情态动词 \textit{skal}）
\end{enumerate}

\textbf{练习 C 答案}
\begin{enumerate}
  \item \textit{Undskyld, er denne plads ledig?} / \textit{Undskyld, er pladsen ledig?}
  \item \textit{Jeg tror, du sidder i den forkerte vogn.}
  \item \textit{Det gør ikke noget, hvor er du sød!}
  \item \textit{Skal du ikke stå af i Aarhus?}
  \item \textit{Jo, det skal jeg!}
\end{enumerate}

\textbf{Før du læser 参考回答 \& 口语复述范文（Resumé）}
\begin{itemize}
  \item \textbf{问题 1}：\textit{Nej, det gør han ikke. Jonas sidder i vogn 7, men hans billet er til vogn 6.}
  \item \textbf{问题 2}：\textit{Emma kommer fra Kina, og hun bor i København. Hun skal stå af på Aarhus Station.}
  \item \textbf{口语复述范文}：\textit{Emma har en billet til vogn 7, plads 14. Men Jonas sidder på hendes plads, fordi han tror, at det er vogn 6. Pladsen ved siden af er ledig, så Jonas bliver siddende og snakker med Emma. Emma kommer fra Kina, og både Emma og Jonas skal stå af i Aarhus.}
\end{itemize}
\end{answerbox}
